\section{复杂度与误差分析}
\label{sec:analysis}

本节给出 DCC-KV 的通信量与计算复杂度，并将端到端近似误差分解为可独立界定的若干项。分析沿
用第~\ref{sec:problem} 节的记号：$N$ 为设备数，$L_s$ 为源块长度，$d_h$、$d_v$ 分别为
KV Head 的 Key 与 Value 维度，$M$ 为代表 Query 数，$d_p$ 为随机投影维度，$B_{s,r}$ 为边
$(s,r)$ 的压缩预算。所有结论针对单个 Transformer 层、单个 KV Head；多层多头情形按层与头
线性叠加。

\subsection{通信量分析}
\label{sec:analysis-comm}

\subsubsection{消息长度}

DCC-KV 的通信由两段组成。第一段是代表 Query 的环形交换：每轮转发一个
$\widehat{Q}_r \in \mathbb{R}^{M \times d_h}$，进行 $N-1$ 轮，故单设备发送量为
\begin{equation}
\label{eq:comm-query}
V_{\text{query}} = (N-1) \, M \, d_h .
\end{equation}
第二段是变长紧凑 KV 的 All-to-Allv：对每条有效边 $(s,r)$，消息长度为
$B_{s,r}(d_h + d_v) + B_{s,r}$，其中最后一项对应质量偏置 $\beta^{s \to r}$。若以
$|\mathcal{E}_s|$ 记源设备 $s$ 的出边数，则其单层发送总量为
\begin{equation}
\label{eq:comm-kv}
V_{\text{kv}} = \sum_{r \in \mathcal{E}_s}
B_{s,r}\,(d_h + d_v + 1).
\end{equation}

作为对照，精确序列并行下源设备 $s$ 必须向同一组目的端发送完整 KV 块，发送量为
\begin{equation}
\label{eq:comm-dense}
V_{\text{dense}} = |\mathcal{E}_s| \, L_s \, (d_h + d_v) .
\end{equation}

若各边预算相同，记为 $B$，则当 $B \ll L_s$ 且 $N$ 不太大时有
\begin{equation}
\label{eq:comm-ratio}
\frac{V_{\text{query}} + V_{\text{kv}}}{V_{\text{dense}}}
\approx \frac{B}{L_s} + \frac{(N-1) M d_h}{|\mathcal{E}_s| L_s (d_h+d_v)} .
\end{equation}
即压缩比主要由 $B/L_s$ 决定，代表 Query 的开销随设备数线性增长但与上下文长度无关，
在长上下文区间可忽略。

\begin{figure}[t]
\centering
\includegraphics[width=\linewidth]{fig4_comm_scaling.pdf}
\caption{单边通信量的解析对比（$d_h+d_v=256$，$B=64$）。精确序列并行的通信量随源块长度
$L_s$ 线性增长；DCC-KV 的紧凑 KV 传输量 $O(B_{s,r}(d_h+d_v))$ 与 $L_s$ 无关，二者相差
$L_s/B$ 倍。代表 Query 的摘要开销 $O(NMd_h)$ 随设备数增长，但与上下文长度无关。
\emph{本图为解析结果，非实测数据。}}
\label{fig:comm}
\end{figure}

\subsubsection{与集合通信实现的匹配}

式~\eqref{eq:comm-query} 的形状规整（所有 $\widehat{Q}_r$ 同形），可直接由 $N-1$ 轮
Ring All-Gather 实现。式~\eqref{eq:comm-kv} 的关键性质是各边长度不同，标准 All-to-All
要求每个接收端预知各来源的消息长度。4.4 节的两阶段实现（先交换长度与缓冲区偏移，再提交
非阻塞收发）把这一约束转化为一次 $O(N^2)$ 的元数据交换，该交换与消息体传输解耦，因此不
进入主通信路径。

\subsection{计算复杂度分析}
\label{sec:analysis-compute}

\subsubsection{代表 Query 选择}

按式~\eqref{eq:proj}--\eqref{eq:repr-q}：随机投影为 $O(L_r d_h d_p)$；$M$ 轮最远点采样，每轮维护并更新当前子集
到所有候选点的最小距离，代价为 $O(M L_r d_p)$。故选择开销为
\begin{equation}
\label{eq:cost-fps}
C_{\text{select}} = O\!\left(L_r d_h d_p + M L_r d_p\right),
\end{equation}
与 $L_r$ 成线性、与 $B_{s,r}$ 无关。注意该开销在所有设备上并行发生。

\subsubsection{紧凑 KV 构造}

对每个目的端 $r$，源设备 $s$ 需计算响应矩阵
$A^{s\to r} = \mathrm{softmax}(\widehat{Q}_r K_s^\top / \sqrt{d_h}) \in
\mathbb{R}^{M \times L_s}$，代价 $O(M L_s d_h)$；RMS 聚合与 Top-$B$ 选择为
$O(M L_s)$ 与 $O(L_s \log B)$。质量偏置拟合的求解规模为 $M \times B$，
投影梯度下降的每次迭代代价 $O(M B)$；Value 回归为 $M \times B$ 的岭回归，
闭式解的构造成本 $O(M B^2 + M B d_v)$。合计
\begin{equation}
\label{eq:cost-build}
C_{\text{build}} = O\!\left(M L_s d_h + M B_{s,r}(d_h + d_v) + M B_{s,r}^2\right).
\end{equation}
式~\eqref{eq:cost-build} 的第一项是主导项。这是目的端条件化的直接代价：源设备 $s$ 需为其
每个出边重复一次构造，总构造代价乘以 $|\mathcal{E}_s|$。

\subsubsection{目的端注意力计算}

目的端 $r$ 的本地块注意力代价为 $O(L_r^2 d_h)$（因果情形减半）。对每个到达的远程紧凑块，
计算式~\eqref{eq:logits} 的 Logit 与式~\eqref{eq:state} 的状态，代价为 $O(L_r B_{s,r}(d_h + d_v))$。全部远程块合计
\begin{equation}
\label{eq:cost-attn}
C_{\text{attn}} = O\!\left(L_r^2 d_h + L_r (d_h + d_v) \textstyle\sum_{s \in \mathcal{E}_r} B_{s,r}\right).
\end{equation}
与精确远程注意力 $O(L_r \sum_{s} L_s (d_h+d_v))$ 相比，式~\eqref{eq:cost-attn} 的第二项
关于 $\sum_s B_{s,r}$ 而非 $\sum_s L_s$，这是压缩在计算侧的收益。

\subsubsection{总览与权衡}

表~\ref{tab:complexity} 汇总各阶段复杂度。需要强调的是，DCC-KV 并不降低本地块的注意力
代价，也不减少层数或头数；其收益集中在跨设备通信量与远程注意力计算量两处，同时以
$|\mathcal{E}_s|$ 倍的构造代价与 $M$ 倍的代表 Query 通信作为代价。当 $B \ll L_s$
且 $M \ll L_r$ 时，这一交换在长上下文区间是划算的。

\begin{table}[t]
\centering
\small
\caption{DCC-KV 各阶段的单层单头复杂度。$B \equiv B_{s,r}$，$L \equiv L_s \approx L_r$。}
\label{tab:complexity}
\begin{tabular}{lll}
\toprule
阶段 & DCC-KV & 精确序列并行 \\
\midrule
代表 Query 交换 & $O(N M d_h)$ & --- \\
紧凑 KV 构造 & $O(|\mathcal{E}_s|(M L d_h + M B^2))$ & --- \\
跨设备通信量 & $O(|\mathcal{E}_s| B (d_h+d_v))$ & $O(|\mathcal{E}_s| L (d_h+d_v))$ \\
远程注意力 & $O(L \sum_s B)$ & $O(L \sum_s L_s)$ \\
\bottomrule
\end{tabular}
\end{table}

\subsection{误差分解与误差界}
\label{sec:analysis-error}

设目的端 $r$ 上某一 Query $q$ 的精确全局注意力为 $\mathrm{Attn}(q)$，DCC-KV 的输出为
$\widehat{\mathrm{Attn}}(q)$。按式~\eqref{eq:global-attn} 的块分解，误差来自三个可分离的环节。以下均针对
单个 Query，且假设所有有效块恰好被处理一次。

\subsubsection{误差的三个来源}

\paragraph{(i) Query 摘要误差。}
源设备仅看到 $\widehat{Q}_r$ 而非完整 $Q_r$。若 $q$ 在投影空间中与所选代表 Query 距离
较远，则据此构造的紧凑 KV 未必适配 $q$。该误差由代表 Query 集合对 $Q_r$ 分布的覆盖度
决定，记作 $\epsilon_{\text{repr}}$。

\paragraph{(ii) 块级质量匹配误差。}
式~\eqref{eq:beta-fit} 在非负约束与 $\lambda_\beta$ 正则下拟合 $m_a^{s\to r}$，一般不能精确复现。定义
\begin{equation}
\label{eq:err-mass}
\epsilon_{\text{mass}}(q) = \sum_{s \in \mathcal{E}_r}
\left| \widehat{M}_s(q) - M_s(q) \right| ,
\end{equation}
其中 $\widehat{M}_s(q) = \sum_{j=1}^{B_{s,r}} \exp(g_{s\to r,j})$。该误差直接改变远程块在
全局 Softmax 分母中的相对权重，是所有误差项中对最终输出影响最直接的一项。

\paragraph{(iii) 条件输出误差。}
式~\eqref{eq:value-reg} 以 $\lambda_v$ 岭回归拟合 $Y^{s \to r}$，得到的是代表 Query 上的最小二乘解，
不保证对任意 $q$ 成立。定义
\begin{equation}
\label{eq:err-out}
\epsilon_{\text{out}}(q) = \sum_{s \in \mathcal{E}_r}
M_s(q) \left\| \widehat{O}_s(q) - O_s(q) \right\|_2 .
\end{equation}

\subsubsection{误差界}

将式~\eqref{eq:global-attn} 的分子分母分别扰动。设
$\widehat{M}_s = M_s + \delta M_s$、$\widehat{O}_s = O_s + \delta O_s$，且记
$S(q) = \sum_{s \in \mathcal{B}_r} M_s(q)$。当 $\sum_s |\delta M_s| < S/2$ 时，
\begin{equation}
\label{eq:err-bound}
\left\| \widehat{\mathrm{Attn}}(q) - \mathrm{Attn}(q) \right\|_2
\;\le\;
\frac{\epsilon_{\text{out}}(q)}{S(q)}
\;+\;
\frac{2\, \epsilon_{\text{mass}}(q)}{S(q)} \cdot \max_{s \in \mathcal{B}_r} \|O_s(q)\|_2 .
\end{equation}
式~\eqref{eq:err-bound} 表明：质量误差通过"相对权重被放大"的机制进入输出，且放大系数是
块输出范数的量级；输出误差则以与自身质量成比例的权重进入。二者均以全局 Softmax 分母
$S(q)$ 归一化，因此远程块占比越大，单块误差的绝对影响越小。

\textbf{该界的首次数值检验。}式~\eqref{eq:err-bound} 此前只被陈述、从未被检验，
而 \S\ref{sec:analysis-error} 的性质 3 又用它论证 DCC-KV 在设备数较大时仍可用。
E13 在「一个压缩源块 $+$ 一个精确上下文块」的两块归并设定下逐 Query 计算两侧：
共 $5760$ 个 Query（$5$ 个种子 $\times$ 关注强度 $\{0,8\}$ $\times$
$B\in\{8,\dots,256\}$ $\times$ $\beta\in\{\text{none},\text{full}\}$），
\textbf{零个反例}，紧致度 $\mathrm{RHS}/\mathrm{LHS}\in[1.00,3.51]$、
中位 $2.10$。判据带 $10^{-9}$ 的相对容差，用于吸收两条数值路径的舍入：
当 $\delta M\equiv0$ 时该界取\textbf{等号}
（$\mathrm{RHS}=\epsilon_{\text{out}}/S=\mathrm{LHS}$），但左侧经逐块 online 归并、
右侧由解析式算出，故纯浮点意义上的「$<$」共 $189$ 例、全部落在取等号的
$10$ 个格子里。若不加该容差，就会把「界恰好取等号」误报成「界被推翻」。

前置条件 $\sum_s|\delta M_s|<S/2$ 的违反率则\textbf{不可忽略}，且其分布本身是结论：
$\beta$ 全量时最高 $16.7\%$，而\emph{不加} $\beta$ 时最高 $100\%$，
并沿预算单调下降（$B=8$ 为 $91.0\%$，$16$ 为 $70.6\%$，$32$ 为 $48.5\%$，
$64$ 为 $15.2\%$，$B\ge128$ 为 $0$）。即\emph{界在误差最大的区间
（小预算且无质量补偿）本就没有保证}，这与该区间的误差主要由 $\beta$ 的缺失
主导完全一致。$\beta$ 的作用因此不止于降低 $\epsilon_{\text{mass}}$：
它同时把该界的适用范围从小预算侧扩展开来。

\begin{figure}[t]
\centering
\includegraphics[width=\linewidth]{fig3_error_decomposition.pdf}
\caption{端到端误差的来源与误差界。质量匹配误差 $\epsilon_{\text{mass}}$ 与输出匹配误差
$\epsilon_{\text{out}}$ 分别作用于块级分母与条件输出，二者均以全局 Softmax 分母 $S(q)$
归一化，故不随设备数发散。代表 Query 覆盖误差 $\epsilon_{\text{repr}}$ 不显式出现在界中，
而是通过前两项间接体现。}
\label{fig:err}
\end{figure}

式~\eqref{eq:err-bound} 中不含 $\epsilon_{\text{repr}}$ 的显式项，因为 Query 摘要误差
通过 $\epsilon_{\text{mass}}$ 与 $\epsilon_{\text{out}}$ 间接体现：$q$ 偏离代表 Query
越远，源设备为它构造的 $(\widehat{M}_s, \widehat{O}_s)$ 就越不准确。这一耦合意味着
$M$ 的选取需要在"摘要开销小"与"$\epsilon_{\text{mass}}, \epsilon_{\text{out}}$ 低"
之间权衡，且该权衡是可测量的（见第~\ref{sec:experiment} 节 E2/E5）。

\subsubsection{若干可陈述的性质}

\paragraph{性质 1（质量守恒的充分条件）。}
若对某条边 $(s,r)$ 有 $B_{s,r} = L_s$，则式~\eqref{eq:key-select} 的 Top-$B$ 选择保留全部 Key，式~\eqref{eq:ck}
退化为 $C_k^{s \to r} = K_s$；此时式~\eqref{eq:beta-fit} 取 $w = \mathbf{1}$、$\beta = \mathbf{0}$ 是
零残差解，式~\eqref{eq:value-reg} 取 $C_v^{s \to r} = V_s$ 亦为零残差解。故 DCC-KV 在预算不裁剪时精确
退化为全注意力，即该框架是精确注意力的一个严格泛化。

需要强调该命题是\textbf{存在性}的充分条件，而非「按默认拟合结果实现即精确」：
预算不裁剪时式~\eqref{eq:value-reg} 的岭回归解 $\widehat{C}_v$ 一般\emph{不等于}
$V_s$，输出侧仍有余项。E2 在 $B=L_s=256$ 处实测：质量侧相对误差为 $0$
（不加 $\beta$）与 $5\times10^{-3}$（$\beta$ 全量），而输出侧相对误差为
$0.070$--$0.225$。也就是说\emph{预算不裁剪时主导误差来自 Value 通路}，
而非选键或 $\beta$——这与 E2 在 $B<L_s$ 处的误差分解定性一致，
也说明「精确退化」只对质量侧无条件成立。

\paragraph{性质 2（无偏置下的质量下界）。}
由于 $w \ge 0$ 且式~\eqref{eq:key-select} 按 RMS 响应选择 Key，被保留 Key 的响应在代表 Query 上系统性地
偏高。故在无偏置情形（$\beta = \mathbf{0}$）下通常有 $\widehat{M}_s(q) \le M_s(q)$，
即压缩后远程块的全局权重被系统性低估，而非随机波动。这为 4.3.2 节引入质量补偿偏置提供
了直接的动机：$\beta$ 需要做的正是把被裁掉的指数响应重新分配给保留的 Key。

\paragraph{性质 3（误差不随块数累积发散）。}
式~\eqref{eq:err-bound} 的右侧以 $S(q)$ 归一化，而 $S(q)$ 本身随有效块数增长。因此，即使
每条边都引入非零误差，只要各边的质量误差不同向叠加，总输出误差不会随设备数 $N$ 线性
放大。这一点是 DCC-KV 在 $N$ 较大时仍可用的必要条件。需要说明的是，该性质依赖误差的
符号分布；若所有边系统性低估质量（如性质 2 所述），则需靠 $\beta$ 拟合抵消，这也解释了
4.3.2 节中正则项选择 $w = \mathbf{1}$ 而非 $w = \mathbf{0}$ 作为收缩目标的设计理由。

\textbf{机制归因的一处更正。}上文把「不随块数放大」归给"各边误差不同向叠加"。
E13 在「每块长度固定、块数 $N$ 增长」的设定下实测（每块 $64$ token、
每块预算 $16$，故总长随 $N$ 增长而每块的压缩强度不变）：
$N$ 由 $1$ 增至 $8$ 时输出侧偏差\emph{不增}（$\beta$ 全量、关注强度 $8$ 下
$0.208\to0.128$；不加 $\beta$、强度 $0$ 下 $0.322\to0.360$，仅 $1.12$ 倍），
而同期 $\sum_s|\delta M_s|/S$ 单调\emph{上升}（$0.038\to0.078$ 与 $0.250\to0.580$），
且 $|\sum_s\delta M_s|\,/\sum_s|\delta M_s|$ 在不加 $\beta$ 时\textbf{恒为} $1.000$——
即各块误差\textbf{完全同向叠加}，一次抵消都没有发生。故真正阻止发散的是
\textbf{分母 $S(q)$ 随有效块数增长}这一项；符号抵消只是在此之上的一层\emph{附加}
裕度（$\beta$ 全量、$N=8$ 时该比值降至 $0.576$），是充分条件而非必要条件。
上文把符号抵消写成条件的表述因此\emph{过强}：即便各边误差同向叠加，
只要 $S(q)$ 同步增长，结论依然成立。

\subsection{异步流水的理论加速上限}
\label{sec:analysis-async}

设单条边的紧凑 KV 传输耗时为 $T_{\text{comm}}$，其对应的注意力计算耗时为
$T_{\text{comp}}$。在纯串行实现（等全部消息到达后统一计算）下，$N-1$ 轮的总耗时约为
\begin{equation}
\label{eq:serial-time}
T_{\text{serial}} \approx (N-1)\left( T_{\text{comm}} + T_{\text{comp}} \right).
\end{equation}
在理想流水下（通信流与计算流完全独立，且任意时刻均有可计算的已到达块），总耗时受两阶段
中较慢者支配：
\begin{equation}
\label{eq:pipeline-time}
T_{\text{async}} \approx \max \Big( (N-1) T_{\text{comm}},\;(N-1) T_{\text{comp}} \Big)
+ T_{\text{comp}} .
\end{equation}
因此理想加速比上界为
\begin{equation}
\label{eq:speedup-bound}
\frac{T_{\text{serial}}}{T_{\text{async}}}
\;\lesssim\;
\frac{T_{\text{comm}} + T_{\text{comp}}}{\max(T_{\text{comm}}, T_{\text{comp}})}
\;=\; 1 + \frac{\min(T_{\text{comm}}, T_{\text{comp}})}{\max(T_{\text{comm}}, T_{\text{comp}})} .
\end{equation}
当 $T_{\text{comm}} \approx T_{\text{comp}}$ 时上界趋近 $2\times$；当两者严重失衡时加速
趋于 $1$。该结论给出一个重要的实践含义：DCC-KV 的异步设计在通信与计算预算相当的区间收益
最大，而在任一侧主导的区间收益有限。第~\ref{sec:experiment} 节 E6 的加速比测量应显式报告
$T_{\text{comm}} / T_{\text{comp}}$ 的实测值，以便判断实测加速偏离上界的原因是流水不充分
还是预算本身失衡。

需要指出，式~\eqref{eq:speedup-bound} 是理论上界，不含通信启动延迟、元数据交换开销、
负载不均衡以及调度抖动。此外，式~\eqref{eq:comm-assoc} 保证的只是数值结果与接收顺序无关，而在线 Softmax
状态归并本身引入浮点舍入差异；该差异的量级已由 E0 界定（\S\ref{sec:exp-cpu}）：
$\mathrm{FP32}$ 下不超过 $7.9\times10^{-7}$、$\mathrm{FP64}$ 下处于机器舍入量级，
且不随乱序试验次数增长；低精度（$\mathrm{FP16}$/$\mathrm{BF16}$）下的失效边界仍需 E8。
因此，异步实现的
实测加速比应低于式~\eqref{eq:speedup-bound}，且二者之差可用于衡量实现质量。

\subsection{小结}

本节确立三点。第一，DCC-KV 把每条边的通信量从 $O(L_s(d_h+d_v))$ 降至
$O(B_{s,r}(d_h+d_v))$，代价是 $|\mathcal{E}_s|$ 倍的构造开销与 $O(NM d_h)$ 的摘要开销，
二者在长上下文区间均可接受。第二，误差可分解为质量误差与输出误差两项并以全局 Softmax
分母归一化，因此不随设备数发散；质量误差通过相对权重的放大机制进入输出。第三，异步流水
的理想加速上界由 $T_{\text{comm}}$ 与 $T_{\text{comp}}$ 的比值决定，该比值是可测的，因而
加速比与上界的差距可以归因。

需要明确的是，本节给出的误差界依赖 $\lambda_\beta$、$\lambda_v$ 与代表 Query 覆盖度等
量。其\emph{紧致程度}已在合成数据上被测量（见上文 E13 与第~\ref{sec:discussion} 节）：
$\mathrm{RHS}/\mathrm{LHS}$ 中位 $2.10$、$95$ 分位 $3.51$、$5760$ 个 Query 零反例，
故它是量级可用的界；但该比值、以及 $\epsilon_{\text{mass}}$ 与 $\epsilon_{\text{out}}$ 的
经验分布，均由合成分布决定，\textbf{在真实模型与真实长文本上尚未测量}。
第~\ref{sec:experiment} 节给出测量方案，其中可在纯 CPU 环境
下复现的部分已标注，需要 GPU 的部分明确标为待执行。
